\subsection{研究内容}
以1个新疆光伏场站的归一化功率为对象，采用15分钟采样与发布间隔，输出未来15--240分钟、步长15分钟的预测，重点报告1、2、3、4小时结果。主体评价限定于正常并网和日照条件，另报全运行时段表现。第2场站独立验证和日前预测为扩展，分别在取得额外场站授权和可追溯业务数值天气预报（NWP）档案后开展。

三项研究内容及其衔接关系如图~\ref{fig:xj-content-framework} 所示。
\begin{figure}[!ht]
\centering
% 项目层研究内容图；原生 TikZ 文字与路径，随 main.tex 官方构建。
% 依据 extraTex/02-content.tex 与 03-methods.tex，不引入新研究承诺。
\begingroup
\definecolor{XJFrameworkBlue}{HTML}{235D82}
\definecolor{XJFrameworkTeal}{HTML}{28777C}
\definecolor{XJFrameworkSlate}{HTML}{48677E}
\definecolor{XJFrameworkMuted}{HTML}{505B65}
\begin{tikzpicture}[
  x=1bp,y=-1bp,
  every node/.style={inner sep=0pt,outer sep=0pt,align=center,
    font={\XJHei\fontsize{9.2bp}{13bp}\selectfont}},
  flow/.style={->,>=stealth,line width=1bp,draw=XJFrameworkBlue}]
  \path[use as bounding box] (0,0) rectangle (414,294);

  \fill[XJFrameworkBlue!7] (0,0) rectangle (414,37);
  \node[text=XJFrameworkBlue,font={\XJHei\bfseries\fontsize{10bp}{14bp}\selectfont}]
    at (207,12) {可用气象信息如何改善概率预测，并传递到爬坡预警？};
  \node[text=XJFrameworkMuted,font={\XJHei\fontsize{8.8bp}{12bp}\selectfont}]
    at (207,27) {新疆典型光伏场站\quad 多源气象信息\quad 短时功率变化与爬坡风险};

  % 三列使用相同层次：研究内容、核心问题、研究要点、预期输出。
  \foreach \x/\col in {0/XJFrameworkBlue,143/XJFrameworkTeal,286/XJFrameworkSlate}{
    \draw[draw=\col,line width=.7bp,rounded corners=2bp,fill=white]
      (\x,50) rectangle +(128,174);
    \fill[\col!10] (\x+.5,50.5) rectangle +(127,25);
    \draw[draw=\col!35,line width=.45bp] (\x+8,121) -- (\x+120,121);
    \fill[\col!6] (\x+.5,184) rectangle +(127,39.5);
  }
  \node[text=XJFrameworkBlue,font={\XJHei\bfseries\fontsize{10.5bp}{14bp}\selectfont}] at (64,63) {研究内容一};
  \node[text=XJFrameworkTeal,font={\XJHei\bfseries\fontsize{10.5bp}{14bp}\selectfont}] at (207,63) {研究内容二};
  \node[text=XJFrameworkSlate,font={\XJHei\bfseries\fontsize{10.5bp}{14bp}\selectfont}] at (350,63) {研究内容三};

  \node[font={\XJHei\bfseries\fontsize{10bp}{14bp}\selectfont}] at (64,96) {多源气象与功率变化\\的跨尺度关系};
  \node[font={\XJHei\bfseries\fontsize{10bp}{14bp}\selectfont}] at (207,96) {天气变化下的\\概率功率预测与校准};
  \node[font={\XJHei\bfseries\fontsize{10bp}{14bp}\selectfont}] at (350,96) {同边际爬坡概率\\估计与预警评价};

  \node[text=XJFrameworkBlue] at (64,135) {哪些信息与尺度有增益？};
  \node[text=XJFrameworkTeal] at (207,135) {如何兼顾覆盖与效率？};
  \node[text=XJFrameworkSlate] at (350,135) {时间依赖有何贡献？};
  \node at (64,164) {可用信息核验\\变量与尺度消融};
  \node at (207,164) {多尺度分位数预测\\静态与滚动校准};
  \node at (350,164) {三组依赖结构对照\\爬坡概率与告警效用};

  \node at (64,204) {输入与尺度选择依据\\信息增益的适用条件};
  \node at (207,204) {原始分位网格、完整误差\\区间可靠性并行诊断};
  \node at (350,204) {事件概率与预警评价\\适用边界与离线原型};
  \draw[flow] (130,98) -- (141,98);
  \draw[flow] (273,98) -- (284,98);

  \node[text=XJFrameworkBlue,font={\XJHei\fontsize{9bp}{13bp}\selectfont}]
    at (207,243) {可用输入\ $\longrightarrow$\ 原始分位网格与完整历史误差\ $\longrightarrow$\ 爬坡概率};
  \node[text=XJFrameworkMuted,font={\XJHei\fontsize{8.8bp}{12bp}\selectfont}]
    at (207,260) {校准区间仅用于并行诊断，不视为完整预测分布已校准。};
  \draw[XJFrameworkBlue!30,line width=.5bp] (0,273) -- (414,273);
  \node[text=XJFrameworkMuted,font={\XJHei\fontsize{8.8bp}{12bp}\selectfont}]
    at (207,286) {拟开展研究；数据按授权获取；输入与标签遵循真实到达时间。};
\end{tikzpicture}
\endgroup

\caption{项目研究内容框架（拟开展研究）}
\label{fig:xj-content-framework}
\end{figure}

\subsubsection{研究内容一：多源气象信息与功率变化的跨尺度关系}
以功率历史与太阳几何为最小输入，关联观测、气象预报和运行状态，逐类加入发布时刻可得的地面观测、NWP及卫星信息。比较各气象源在慢变与快速扰动中的增益及其天气、时距差异，以同预算变量和尺度消融控制容量影响，形成供内容二选取输入的数据质量说明与尺度依据。

\subsubsection{研究内容二：天气变化下的概率功率预测与校准}
以内容一保留的输入建立多尺度分位数模型，输出中位预测和功率分位数。在独立时段选模，比较静态与滚动校准在高指数平稳、辐照快速变化过程中的覆盖偏差、区间宽度与概率评分，检查多数天气改善或区间过宽是否掩盖失配；云和沙尘过程据外部证据另述。输出逐时距概率预测、校准诊断和完整历史误差向量，供内容三使用。

\subsubsection{研究内容三：保持边际一致的爬坡概率估计与预警评价}
各时距使用相同分位数样本，仅改变跨时距排序，构造独立排序、总体历史秩和天气条件历史秩三组轨迹。以功率净变化分布比较时间依赖及天气条件选择的贡献，并以轻量直接增量模型对照。告警阈值在验证集确定，分别评价概率质量和告警效用，形成可回放天气过程、功率分布与预警结果的离线原型。

\subsection{研究目标}
目标是建立适用于新疆场站的短时概率预测与爬坡分析方法，明确适用天气与失效条件。通过同输入时限、同预算的跨时段消融检验气象源和尺度相对功率历史的稳定增益；以留出时段、天气分组的配对比较量化滚动相对静态校准的覆盖与效率；以同边际排序对照、直接增量基线和事件级不确定性分析检验依赖对事件概率及提前识别的贡献。

成果包括数据清单、模型与评价程序、案例报告及离线原型。以可重复改善判断有效性，不预设精度增幅；报告各时距、天气的阴性结果和适用边界，地域化结论以新疆实测为据。

\subsection{拟解决的关键科学问题}
第一个问题是：\textbf{在发布时刻可用的信息集合固定后，气象变量的预测增益怎样随时间尺度、天气状态和预报时距变化？}以同输入时限、同预算消融检验可泛化增益；相关和滞后关系不作物理因果解释。

第二个问题是：\textbf{天气转换造成误差分布变化时，利用已经到达的历史标签进行校准，能否兼顾覆盖水平和区间效率？}检验标签到达时限下的适应延迟和小样本影响，单列少见天气证据，不由总体覆盖推断各类天气均可靠。

第三个问题是：\textbf{边际概率预测保持一致时，时间依赖对爬坡事件概率可靠性的贡献及其失效条件是什么？}以同边际轨迹的事件评分差异识别依赖贡献；差异不稳定时采用更简单的事件预测路线。
